\documentclass[10pt,a4paper]{amsart}
\usepackage{amsmath,amssymb,amsthm,graphicx}
% Retained minimal preamble for the Stacks Project chapter more-algebra.tex (More on Algebra).
% The chapter file pinned-original-source.tex carries only \input{preamble} on its first line, so no
% preamble text lives inside the chapter itself.  The Stacks Project uses one single shared preamble for
% every chapter; the snapshot stages it identically as sources/stacks-chow/preamble.tex and
% sources/stacks-etale-cohomology/preamble.tex.  The block below is copied mechanically, line for line,
% from sources/stacks-chow/preamble.tex lines 14-18 (the Xy-pic setup) and 154-217 (the native theorem
% declarations together with the mathematical macros).  Nothing was added, removed or redefined; the
% wrapper class and the amsmath/amssymb/amsthm/graphicx base packages are supplied by the assembler.
% Chapter: more-algebra.tex -- More on Algebra.
% For commutative diagrams we use Xy-pic
\usepackage[all]{xy}

% We use 2cell for 2-commutative diagrams.
\xyoption{2cell}

% Theorem environments.
%
\theoremstyle{plain}
\newtheorem{theorem}[subsection]{Theorem}
\newtheorem{proposition}[subsection]{Proposition}
\newtheorem{lemma}[subsection]{Lemma}

\theoremstyle{definition}
\newtheorem{definition}[subsection]{Definition}
\newtheorem{example}[subsection]{Example}
\newtheorem{exercise}[subsection]{Exercise}
\newtheorem{situation}[subsection]{Situation}

\theoremstyle{remark}
\newtheorem{remark}[subsection]{Remark}
\newtheorem{remarks}[subsection]{Remarks}

\numberwithin{equation}{subsection}

% Macros
%
\def\lim{\mathop{\mathrm{lim}}\nolimits}
\def\colim{\mathop{\mathrm{colim}}\nolimits}
\def\Spec{\mathop{\mathrm{Spec}}}
\def\Hom{\mathop{\mathrm{Hom}}\nolimits}
\def\Ext{\mathop{\mathrm{Ext}}\nolimits}
\def\SheafHom{\mathop{\mathcal{H}\!\mathit{om}}\nolimits}
\def\SheafExt{\mathop{\mathcal{E}\!\mathit{xt}}\nolimits}
\def\Sch{\mathit{Sch}}
\def\Mor{\mathop{\mathrm{Mor}}\nolimits}
\def\Ob{\mathop{\mathrm{Ob}}\nolimits}
\def\Sh{\mathop{\mathit{Sh}}\nolimits}
\def\NL{\mathop{N\!L}\nolimits}
\def\CH{\mathop{\mathrm{CH}}\nolimits}
\def\proetale{{pro\text{-}\acute{e}tale}}
\def\etale{{\acute{e}tale}}
\def\QCoh{\mathit{QCoh}}
\def\Ker{\mathop{\mathrm{Ker}}}
\def\Im{\mathop{\mathrm{Im}}}
\def\Coker{\mathop{\mathrm{Coker}}}
\def\Coim{\mathop{\mathrm{Coim}}}

% Boxtimes
%
\DeclareMathSymbol{\boxtimes}{\mathbin}{AMSa}{"02}

%
% Macros for moduli stacks/spaces
%
\def\QCohstack{\mathcal{QC}\!\mathit{oh}}
\def\Cohstack{\mathcal{C}\!\mathit{oh}}
\def\Spacesstack{\mathcal{S}\!\mathit{paces}}
\def\Quotfunctor{\mathrm{Quot}}
\def\Hilbfunctor{\mathrm{Hilb}}
\def\Curvesstack{\mathcal{C}\!\mathit{urves}}
\def\Polarizedstack{\mathcal{P}\!\mathit{olarized}}
\def\Complexesstack{\mathcal{C}\!\mathit{omplexes}}
% \Pic is the operator that assigns to X its picard group, usage \Pic(X)
% \Picardstack_{X/B} denotes the Picard stack of X over B
% \Picardfunctor_{X/B} denotes the Picard functor of X over B
\def\Pic{\mathop{\mathrm{Pic}}\nolimits}
\def\Picardstack{\mathcal{P}\!\mathit{ic}}
\def\Picardfunctor{\mathrm{Pic}}
\def\Deformationcategory{\mathcal{D}\!\mathit{ef}}

% --- appended 2026-10-12 (Chief): the project own definitions of the reference/slogan/history
% display environments, copied mechanically from the shared Stacks preamble lines 8-12
% (\usepackage{verbatim} plus the three \newenvironment lines, which the chapter preamble keeps
% verbatim from the project preamble). Nothing added, removed or redefined: these environments
% carry no mathematical content and the project itself typesets them as comments. Without these
% three lines a counted statement whose own text uses them fails with Environment slogan undefined.
\usepackage{verbatim}
\newenvironment{reference}{\comment}{\endcomment}
\newenvironment{slogan}{\comment}{\endcomment}
\newenvironment{history}{\comment}{\endcomment}
\usepackage{cleveref}
\crefname{originallemma}{original lemma}{original lemmas}
\Crefname{originallemma}{Original lemma}{Original lemmas}
\makeatletter
\newlabel{algebra-lemma-finite-gl-dim-finite-dim-regular}{{00OE}{source}{}{sourcepointer-1}{}}
\newlabel{algebra-lemma-finite-gl-dim-finite-dim-regular@cref}{{[originallemma][0][]00OE}{[1][0][]source}{}{}{}}
\makeatother
\crefname{originallemma}{original lemma}{original lemmas}
\Crefname{originallemma}{Original lemma}{Original lemmas}
\makeatletter
\newlabel{derived-lemma-morphisms-from-projective-complex}{{064B}{source}{}{sourcepointer-2}{}}
\newlabel{derived-lemma-morphisms-from-projective-complex@cref}{{[originallemma][0][]064B}{[1][0][]source}{}{}{}}
\makeatother
\title{There are no maps between bounded above complexes over a regular ring when their support ranges are separated by the dimension}
\author{The Stacks Project Authors}
\date{stacks-project:more-algebra.tex:local-snapshot-20261012; GNU FDL 1.2 or later}
\begin{document}
\maketitle
\setcounter{subsection}{1}
\noindent\textit{Editorial note.} All delivered mathematical text is the author's own, copied from the pinned Stacks Project chapter file more-algebra.tex; no statement or proof step was rewritten, reordered or removed. The counted statement is the labeled lemma environment at source lines 20766-20772 (canonical tag 0FXJ) and the counted proof is the proof environment at source lines 20774-20797, which immediately follows it and is its own proof.\par\medskip
\noindent\textit{Editorial note.} The printed text uses the letter k both for the vanishing bound of the first complex and inside the phrase there exists an k; the wording is copied exactly as printed.\par\medskip
\noindent\textit{Editorial note.} The preamble span runs to line 90, the current end of the retained preamble, which includes the appended environment definitions; the counted text does not use them, and the span is taken whole so that the file version is unambiguous.\par\medskip
\noindent\textit{Editorial note.} The source URL is the source-level Stacks Project URL, as required when more than one card is produced from one source; the per-claim tag is carried by claim\_key and by this note. The counted statement and proof cite no bibliography entry, so this card needs no bibliography.\par\medskip
\noindent\textit{Allowed standard prerequisites (editorial).}\begin{itemize}
\item A bounded above complex may be replaced by its truncation above the vanishing bound in low degrees without changing the Ext groups in question. used at proof lines 20774-20778. Puts the representing complex in the normal form in which the top degree argument applies.
\item A short exact sequence of complexes with the shift of a term in the middle gives a distinguished triangle, so the vanishing statement can be checked on the two outer pieces. used at proof lines 20778-20783. Reduces the statement by induction on the number of nonzero terms to the case of a single module placed in one degree.
\item Over a regular ring of finite dimension every module has a projective resolution of length at most the dimension, and a map of complexes from a bounded complex of projectives concentrated in a range to a complex vanishing above a lower bound is zero. used at proof lines 20783-20796. Provides the final vanishing, which is exactly the statement for the single module case.
\end{itemize}\medskip
\ifdefined\hypertarget\hypertarget{sourcepointer-1}{}\fi
\noindent\textit{Source reference.} Original lemma 00OE: sources/stacks-chow/tags.txt, source lines 892--892. The reference gives finite global dimension of the regular ring, used at proof lines 20784-20785 to bound the length of the projective resolution; the induction and the final comparison are inside the counted span.\par\medskip
\ifdefined\hypertarget\hypertarget{sourcepointer-2}{}\fi
\noindent\textit{Source reference.} Original lemma 064B: sources/stacks-chow/tags.txt, source lines 7860--7860. The reference identifies maps in the derived category with maps of complexes from a bounded complex of projectives, used at proof lines 20796-20797 to conclude the vanishing from the disjointness of supports; it is the last step and the preparation for it is inside the counted span.\par\medskip
\begin{lemma}
\label{lemma-ext-regular}
Let $R$ be a Noetherian regular ring of dimension $d < \infty$.
Let $K, L \in D^-(R)$. Assume there exists an $k$ such that
$H^i(K) = 0$ for $i \leq k$ and $H^i(L) = 0$ for $i \geq k - d + 1$.
Then $\Hom_{D(R)}(K, L) = 0$.
\end{lemma}

\begin{proof}
Let $K^\bullet$ be a bounded above complex representing $K$, say
$K^i = 0$ for $i \geq n + 1$. After replacing $K^\bullet$ by
$\tau_{\geq k + 1}K^\bullet$ we may assume $K^i = 0$ for $i \leq k$.
Then we may use the distinguished triangle
$$
K^n[-n] \to K^\bullet \to \sigma_{\leq n - 1}K^\bullet
$$
to see it suffices to prove the lemma for $K^n[-n]$ and
$\sigma_{\leq n - 1}K^\bullet$. By induction on $n$, we conclude
that it suffices to prove the lemma in case $K$ is represented by
the complex $M[-m]$ for some $R$-module $M$ and some $m \geq k + 1$.
Since $R$ has global dimension $d$ by Algebra, Lemma
\ref{algebra-lemma-finite-gl-dim-finite-dim-regular}
we see that $M$ has a projective resolution
$0 \to P_d \to \ldots \to P_0 \to M \to 0$.
Then the complex $P^\bullet$ having $P_i$ in degree $m - i$
is a bounded complex of projectives representing $M[-m]$.
On the other hand, we can choose a complex $L^\bullet$ representing
$L$ with $L^i = 0$ for $i \geq k - d  + 1$.
Hence any map of complexes $P^\bullet \to L^\bullet$ is zero.
This implies the lemma by Derived Categories, Lemma
\ref{derived-lemma-morphisms-from-projective-complex}.
\end{proof}

\end{document}
